\documentclass[11pt,a4paper]{article}
\usepackage[margin=2.35cm]{geometry}
\usepackage[T1]{fontenc}
\usepackage[utf8]{inputenc}
\usepackage{lmodern}
\usepackage{microtype}
\usepackage{enumitem}
\usepackage{booktabs}
\usepackage{array}
\usepackage[hidelinks]{hyperref}
\setlength{\parindent}{0pt}
\setlength{\parskip}{5.5pt}
\setlist[itemize]{leftmargin=1.4em,itemsep=2pt,topsep=3pt}
\setlist[enumerate]{leftmargin=1.6em,itemsep=2pt,topsep=3pt}
\newcommand{\closed}[1]{\vspace{4pt}\textbf{#1}\par}
\newcommand{\ruletext}[1]{\begin{quote}\bfseries #1\end{quote}}
\begin{document}

{\LARGE\bfseries DDTA - Working Metamodel - Decision}\par
\vspace{4pt}
{\large\bfseries NON CANONICAL WORKING DRAFT - REVISION 7}\par
\vspace{4pt}
Revision 7 is a layering correction. It keeps the Decision-specific semantic core and the hierarchy/corpus invariants introduced by revision 6, but removes representation, realization, semantic-authority and authoring-review rules from the Decision-specific closed set. Those rules remain part of the study in the separate cross-cutting model-layering foundation document.

\section{Purpose and evidence boundary}
This study asks what a governed \textbf{Decision} means immediately below a Macro Requirement. Construction evidence remains the sixteen historical ThreatForge ADRs of MR-0001/MR-0002, followed by the independent facial-access example. ThreatForge MR-0003, MR-0004 and MR-0005 remain sequential holdouts and are not opened by this revision.

Revision 6 added two genuine cross-document semantic invariants after the independent example exposed ambiguous parentage and redundant Decision contribution. The subsequent ThreatForge regression exposed a different problem in the study itself: a realization principle (executable projections) was being applied as if it were semantic truth of every governed DDTA project. Revision 7 corrects that layering without changing the meaning of Decision.

\section{Semantic position below Macro Requirement}
\begin{verbatim}
Macro Requirement
    major governed concern/result
         |
         v
Decision
    significant chosen narrowing of that concern
         |
         v
Requirement-level semantics (to be studied next)
    independently checkable obligations / downstream effects
\end{verbatim}
The last level is intentionally not frozen here. Requirement remains the working name of the next document family.

\section{Decision-specific semantic rules CLOSED for the construction candidate}
\closed{D-CLOSED-01 - Purpose of Decision}
\ruletext{A Decision records one significant choice that resolves an unresolved problem or tension within one parent Macro Requirement and narrows that macro concern.}
Its subject may be a method, policy, convention, operating practice, semantic choice, representation choice, technology or architecture when that is the actual choice being governed.

\closed{D-CLOSED-02 - Minimal semantic core}
\begin{verbatim}
Decision
|- title
|- context
|- decision
`- consequences
\end{verbatim}
All four Decision-specific semantic attributes are required and non-null in the current candidate. Parentage is an explicit semantic relation to the Macro Requirement rather than duplicated prose.

\closed{D-CLOSED-03 - Decision-local Context}
\ruletext{Context describes the unresolved local decision problem, tension or ambiguity that makes this specific choice necessary inside the already-governed parent concern.}
If Context mostly explains why the parent MR exists, it is too high-level. Context may preserve historically true circumstances of the choice, but it must not become a cumulative copy of inherited knowledge.

\closed{D-CLOSED-04 - One coherent chosen position}
\ruletext{Decision states one coherent chosen position that can meaningfully be accepted, rejected or superseded as a conceptual unit.}
Split diagnostic: if A can change while B remains valid and both have independently meaningful trade-offs, the candidate likely contains more than one Decision.

\closed{D-CLOSED-05 - Consequence awareness}
\ruletext{Consequences record relevant effects and trade-offs accepted because of the choice.}
Benefit, Cost, Risk and Constraint remain useful authoring labels but are not frozen semantic subtypes.

\closed{D-CLOSED-07 - No mandatory embedded Alternatives or Rationale fields}
The evidence does not justify mandatory \texttt{alternatives[]} or a separate \texttt{rationale} semantic field. Substantial alternatives may exist as separately governed candidate Decisions. Reasoning is expected to be recoverable from Context + Decision + Consequences unless a future counterexample demonstrates otherwise.

\section{Cross-document Decision-set invariants CLOSED}
These invariants belong to the MR--Decision relation and to the governed Decision set. They are not additional fields inside a Decision.

\closed{H-CLOSED-01 - Unique semantic ownership}
\ruletext{Each Decision has exactly one Macro Requirement whose concern is the natural semantic owner of the unresolved problem resolved by that Decision.}
Different MRs may mention the same domain objects and vocabulary, and a Decision may be relevant to other concerns. Relevance does not create multiple hierarchical ownership.

Unique-parent test: if a Decision can move under another MR without materially changing its Context or chosen meaning, review whether the Decision is too broad, redundant/derived, or whether MR boundaries fail to discriminate responsibility. A genuinely cross-cutting choice is not resolved by copying it under multiple MRs.

\closed{H-CLOSED-02 - Non-redundant Decision contribution}
\ruletext{Each governed Decision contributes a significant choice that is not already owned by another applicable Decision and is not completely determined by its parent Macro Requirement plus the already governed choices.}
Two Decisions may share terminology, actors, consequences or evidence while owning different choices. The prohibited condition is duplicate semantic ownership of the same choice or promotion of a logically derived statement into a second Decision.

Decision-contribution test: if removing a candidate would not remove a significant project choice that cannot be reconstructed from the parent concern and remaining Decisions, review it as a possible duplicate, derived consequence, explanation or downstream obligation.

\section{Three semantic validation levels}
\begin{verbatim}
1. Local Decision validity
   semantic core / local problem / coherent chosen position

2. Hierarchical coherence
   one natural MR owner / discriminable parentage

3. Corpus coherence
   non-redundant contribution / no duplicate choice ownership
\end{verbatim}
The second and third levels inherently require governed context beyond one file. Tooling may assist review at scale but does not own these semantics.

\section{Rules deliberately relocated from this metamodel}
Revision 7 retains the historical identifiers but changes their owner:
\begin{itemize}
\item \textbf{D-CLOSED-06 - No semantic patching}: authoring/review heuristic, not a Decision invariant. Non-goals remains outside the Decision semantic core.
\item \textbf{D-CLOSED-08 - Lifecycle is not Decision-specific}: cross-cutting common-governance boundary. Lifecycle/status/history will be modeled separately.
\item \textbf{D-CLOSED-09 - Semantic-owner separation}: cross-cutting semantic-authority principle applying to DDTA, governed projects and tools.
\item \textbf{R-CLOSED-01..04}: representation and realization rules; see the model-layering foundation document.
\end{itemize}
Relocation is not rejection. It prevents rules about serialization or tool architecture from being treated as ontological properties of Decision.

\section{Decision-to-downstream boundary - still partly OPEN}
Working rule:
\ruletext{A Decision chooses; downstream governed documents make the chosen consequences operational and independently checkable.}
The exact semantics of the next Requirement level are not yet closed.

\begin{verbatim}
choice needed to explain what was adopted
    -> Decision

independently checkable obligation
    -> candidate next-level semantics

rollout / migration order
    -> planning

separable choice with its own trade-offs
    -> another Decision candidate
\end{verbatim}
Technical vocabulary is not forbidden; the question is whether the detail belongs to the choice or to downstream realization.

\section{Historical integrity}
Decision Context may preserve circumstances true when the choice was taken. Later replacement should normally be represented through the future lifecycle/governance model rather than rewriting history to look current. Temporary rollout sequencing remains planning rather than Decision semantics.

\section{Decision review tests after revision 7}
The following tests assess the conceptual Decision model:
\begin{enumerate}
\item \textbf{Parent}: does the candidate narrow one parent MR?
\item \textbf{Unique parent}: is that MR the single natural semantic owner?
\item \textbf{MR discriminability}: if parentage is ambiguous, is the defect in the candidate or in MR boundaries?
\item \textbf{Contribution}: would removing it remove a significant choice not already governed or derivable?
\item \textbf{Decision-local Context}: does Context explain this unresolved choice rather than the whole MR?
\item \textbf{Explicit choice}: is an adopted position actually present?
\item \textbf{Coherent choice}: can the choice change as one conceptual unit?
\item \textbf{Consequences}: are material effects/trade-offs represented?
\item \textbf{Downstream boundary}: are independently checkable obligations descending rather than bloating the Decision?
\item \textbf{Historical integrity}: does Context preserve why the choice existed without carrying temporary rollout state?
\end{enumerate}
Authoring heuristics, semantic-authority review, representation conformance and tooling checks are intentionally documented outside this Decision-specific test list.

\section{Open Decision questions before candidate closure}
\begin{enumerate}
\item Does re-review of the independent facial-access corpus expose a counterexample to the local core or H-CLOSED-01/H-CLOSED-02?
\item Where exactly is the threshold between precision intrinsic to a choice and next-level independently checkable detail?
\item Are consequence labels semantic categories or only authoring aids?
\item How should legitimate cross-MR relevance eventually be represented without multiple hierarchical ownership?
\end{enumerate}

\section{Validation sequence after layering correction}
\begin{enumerate}
\item preserve the frozen facial-access D1 corpus;
\item preserve revision 7 as a layering checkpoint;
\item re-read the ThreatForge construction ADR regression without treating realization principles as metamodel semantics;
\item write a corrected candidate ADR set only after that regression is stable;
\item regress all prior evidence if Decision semantics themselves change;
\item only then proceed to sequential ThreatForge MR-0003/MR-0004/MR-0005 holdouts.
\end{enumerate}

\section{Explicitly deferred}
\begin{itemize}
\item lifecycle/status value sets and transitions;
\item Requirement-specific semantic model;
\item cross-MR relevance relation, if needed;
\item cross-document persistence/effective-context relation model;
\item concrete representation mapping and serialization closure;
\item concrete ThreatForge support/conformance relation model;
\item concrete similarity/search algorithms and thresholds;
\item ThreatForge migration/refactor.
\end{itemize}

\end{document}
